% Pista alg-23s-q3 · Àlgebra
% Procedència: PAU Catalunya 2023 · Convocatòria de setembre · Sèrie 2 · Qüestió 3
\documentclass[11pt,a4paper]{article}
\usepackage{pista}
\pistainfo{Catalunya 2023}{Convocatòria de setembre}{Sèrie 2}

\begin{document}

\section*{\textcolor{blauPista}{Qüestió 3}}

\noindent Sigui el sistema d'equacions lineals
\[
\left.\begin{array}{r}
2x + y = 1 + z \\
my + z = 2 - x \\
mz + 3 = 3x + y
\end{array}\right\},
\]
en què $m$ és un nombre real.

\subsection*{\textit{a)} \normalfont Discutiu el sistema segons els valors del paràmetre $m$. \hspace*{\fill}\textnormal{[1,25~punts]}}

\pista{Abans de calcular res, passa totes les incògnites a l'esquerra, en l'ordre $x$, $y$, $z$, i els nombres a la dreta. Per exemple, la primera equació, $2x + y = 1 + z$, queda així:
\[
2x + y - z = 1
\]
Fes el mateix amb les altres dues. A continuació, escriu la matriu $A$ i la matriu ampliada $A^{'}$.

\medskip
\textbf{Atenció:} quan un terme passa a l'altre costat, canvia de signe. Vigila-ho sobretot a la tercera equació.

\medskip
Llavors, discuteix el sistema com sempre: cada valor d'$m$ que anul·la el determinant d'$A$ s'estudia a part, i la resta de valors es tracten tots en un sol cas, perquè el sistema és determinat i té solució única.}

\subsection*{\textit{b)} \normalfont Resoleu el sistema, si té solució, per al cas $m = 1$. \hspace*{\fill}\textnormal{[1,25~punts]}}

\pista{En primer lloc, mira a la discussió de l'apartat \textit{a)} en quin cas és $m = 1$. Així sabràs si el sistema té solució, i quantes en té.

\medskip
A continuació, substitueix $m = 1$ al sistema que has reorganitzat. Fixa't en la 1a i la 3a equació: els termes amb $y$ i amb $z$ són iguals o bé oposats. Si són iguals, resta-les; si són oposats, suma-les. En tots dos casos desapareixen la $y$ i la $z$, i obtindràs directament el valor de $x$.

\medskip
Finalment, substitueix aquest valor a la 1a i a la 2a equació: et quedarà un sistema de dues equacions amb $y$ i $z$.}

\end{document}
